\documentclass[10pt,a4paper]{amsart}
\usepackage[margin=19mm]{geometry}
\usepackage{amsmath,amssymb,mathtools}
\usepackage[scr=rsfs,cal=euler]{mathalfa}
\usepackage{xfrac}
\usepackage[unicode,hidelinks,bookmarks=false]{hyperref}
\newcommand{\ie}{i.e.\ }
\newcommand{\cf}{cf.\ }
\newcommand{\resp}{respectively\ }
\newcommand{\Subsection}{\subsection}
\providecommand{\nobreakdash}{\nobreak\hbox{-}}
\setlength{\emergencystretch}{2em}
\multlinegap0pt
\usepackage{bm}
\usepackage{bbm}
\usepackage{dsfont}
\usepackage{xspace}
\usepackage{cancel}
\usepackage{mathtools}
\usepackage{amsthm}
\makeatletter
\@ifundefined{lemma}{\newtheorem{lemma}{Lemma}}{}
\makeatother
\title[Elliptic islands and zero measure escaping orbits in a class]{Elliptic islands and zero measure escaping orbits in a class of outer billiards}
\author{Zaicun Li}
\date{Source: arXiv:2506.18433v1 (CC BY 4.0)}
\begin{document}
\maketitle

\noindent\emph{Source and licence.} Elliptic islands and zero measure escaping orbits in a class of outer billiards. Version arXiv:2506.18433v1 (CC BY 4.0), distributed under the Creative Commons Attribution 4.0 International licence, as recorded in the arXiv licence metadata for this exact version and shown on the abstract page https://arxiv.org/abs/2506.18433v1. Full rights notice: licenses/arxiv-2506.18433-CC-BY-4.0.txt.\par\smallskip

\noindent\textit{Editorial scope.} Counted unit: the Lemma whose source label is \texttt{conj\_conjugacy} in the article (source0.tex lines 533--562, source label \texttt{conj\_conjugacy}), delivered together with the article's own complete original proof of it (source0.tex lines 563--632, one \texttt{proof} environment of 70 lines). No mathematics is written, repaired or re-derived: statement and proof are the article's own text line for line. Required context retained uncounted: formab (lines 484--488). Every definition, display and label of those lines is retained unchanged. Omitted from this package and not redistributed: 1--483, 489--532, 633--2237. Nothing inside a retained excerpt is omitted or altered: every retained body is delivered exactly as the article's lines are written, so no omission receipt of any kind accompanies this package. Citations: the retained excerpts carry no citation command at all, so no bibliography entry of the article is transcribed and no reference is redistributed. Reference adaptation: none was needed and none was made; every reference inside the delivered unit resolves inside it, so no unresolved reference or citation is produced. Printed slips kept literal: no printed word, symbol, hypothesis or number of the counted statement or of its proof was altered. Layout and class: the article's own class is \texttt{article} with packages including caption, booktabs, verbatim, mathtools, hyperref, float, amssymb, graphicx, amsmath, amsthm, amsfonts, tikz, scalefnt, lineno; the wrapper is the lane's pinned \texttt{amsart} editorial wrapper at 10pt on A4, which restates the theorem-like environments the retained text needs and adds the article's own macro definitions that the retained text needs (none), with extra wrapper packages (bm, bbm, dsfont, xspace, cancel, mathtools). The delivered numbering is the unit's own. Assets: no retained span contains a figure environment, a graphics-inclusion command, a PDF-page inclusion, a table or a TikZ picture; no image file is copied into this package, the package declares no asset dependency and no image file is redistributed with it. Licence: version arXiv:2506.18433v1 of this article is distributed under the Creative Commons Attribution 4.0 International licence, as recorded in the arXiv licence metadata for this exact version and shown on the abstract page https://arxiv.org/abs/2506.18433v1. A full rights notice is delivered at licenses/arxiv-2506.18433-CC-BY-4.0.txt. Characters: every retained excerpt is pure ASCII, so the delivered engine, XeLaTeX with its default Latin Modern text font, reports no missing character.
	\begin{align}\label{formab}
		 T_1=H_{2,1}&\circ H_1\circ \mathcal{F} \circ H_1^{-1}\circ H_{2,1}^{-1}: (a,b)\mapsto (a',b'),\notag\\
		a'&=\lambda a + \sum_{j = 2}^{3}{G_{1,j}(a,b)}+\mathcal{O}_4\left( n^{- 3} \right),\notag\\
		b'&=\lambda^{- 1}b + \sum_{j = 2}^{3}{G_{2,j}(a,b)}+\mathcal{ O}_4\left( n^{- 3} \right),
	\end{align}

\begin{lemma}
	There is a conjugacy $H_{2,2}: (a,b)\rightarrow (\xi, \eta)$ given by the following implicit form,
	\begin{equation}\label{conj_conjugacy}
		a = \xi + \sum_{m = 2}^{3}{\phi_{m}(\xi_1,\eta_1)},\ b = \eta + \sum_{m = 2}^{3}{\psi_{m}(\xi_1,\eta_1)},
	\end{equation}
	where \(\phi_{m}\) and \(\psi_{m}\) are homogeneous polynomials of
	degree \(m\), and under this conjugacy, $T_2=H_{2,2}\circ T_1 \circ H_{2,2}^{-1} $ takes the form
	\begin{align}\label{formBirk}
		\begin{split}
			\xi_1'= u\xi_1 + \mathcal{ O}_4\left( n^{- 3} \right),\\
			\eta_1'= v\eta_1 +\mathcal{ O}_4\left( n^{- 3} \right),\\ 
			u = \overline{v} = e^{i(\alpha + \alpha_{2}\xi_1\eta_1)},
		\end{split}
	\end{align}
	where \(e^{i\alpha} = \lambda\).
	
	Moreover, if we let
	\begin{equation}\label{domain of xi}
		V=0.98P\left(	\frac{3}{2048},  -\frac{3}{2048}+\frac{3\sqrt{2}i}{1024}, -\frac{3}{1024},  \frac{3}{2048}-\frac{3\sqrt{2}i}{1042}\right)
	\end{equation}
	be a complex parallelogram, and let $\widetilde{V}=\left\{(\xi,\eta)\in \mathbb{C}^2, \xi=\overline{\eta},\xi\in V\right\}$.
	Then for any $n$ large enough,
	$$H_{2,2}\left(\widetilde{U}\right)\supset \widetilde{V},$$ and we have,
	\begin{equation}\label{complex_conj}
		{\overline{\phi}}_{m}(\eta_1,\xi_1) = \psi_{m}(\xi_1,\eta_1).
	\end{equation}
	When $a$ and $b$ are complex conjugate numbers, $\xi_1$ and $\eta_1$ will also be complex conjugate numbers.
	
	Furthermore, the polynomials $\phi_m$ and $\psi_m$ are of order $\mathcal{O}(n^{-m+1})$. So, this conjugacy is almost identity when $n$ goes to infinity.
\end{lemma}

\begin{proof}
We assume the conjugacy is in the form of (\ref{conj_conjugacy}) and let
\(\phi_{m}(\xi_1,\eta_1) = \sum_{l = 0}^{m}a_{m}^{l}\xi_1^{l}\eta_1^{m - l}\),
\(\psi_{m}(\xi_1,\eta_1) = \sum_{l = 0}^{m}b_{m}^{l}\xi_1^{m - l}\eta_1^{l}\). Further,  we assume that 
\begin{equation}\label{conj_relation}
\overline{a_{m}^{l}} = b_{m}^{l},
\end{equation}
then condition (\ref{complex_conj}) will be satisfied.
Assume that this conjugacy will turn (\ref{formab}) into the form (\ref{formBirk}).

 We then solve the following equation. 
\begin{equation}
	H_{2,2}^{-1}\circ T_2=T_1 \circ H_{2,2}^{-1}.
\end{equation}
To solve this equation, we expand both sides into polynomials of
\((\xi_1,\eta_1)\) to the fourth degree, then the equation involving \(\xi_1\) turns out to be
\begin{equation}\label{Birk_eq}
	\begin{split}
&\lambda\left( \xi_1 + \sum_{m = 2}^{3}{\phi_{m}(\xi_1,\eta_1)} \right) + \sum_{l = 2}^{3}G_{1,l} \left( \xi_1 + \sum_{m = 2}^{3}{\phi_{m}(\xi_1,\eta_1)},\eta_1 + \sum_{m = 2}^{3}{\psi_{m}(\xi_1,\eta_1)} \right)\\
 & = \lambda\left( 1 + \sum_{m=1}^{\infty}{\left(i\alpha_{2}\xi_1\eta_1\right)}^m/m!  \right)\xi_1  \\
 &+ \sum_{m = 2}^{3}{\phi_{m}\left( \lambda\left( 1 + i\alpha_{2}\xi_1\eta_1 \right)\xi_1 +\mathcal{O}_4(n^{-3}), \lambda^{- 1}\left( 1 - i\overline{\alpha_{2}}\xi_1\eta_1 \right)\eta_1 +\mathcal{O}_4(n^{-3}) \right)}\\
 &+\mathcal{O}_4(n^{-3}),
 \end{split}
\end{equation}

It should be pointed out that the equation for \(\eta_1\) will give 
exact the complex conjugate version of (\ref*{Birk_eq}), so there is no need to list
it here. Our goal is to find $a_m^l$ and $b_m^l$ so that all the coefficients of homogeneous polynomials of $\xi_1$ and $\eta_1$ of degree less than four cancel out in (\ref{Birk_eq}). This, together with (\ref{conj_relation}), would lead us to the following equations:
\begin{align}
    {\rm coefficient\ of\ } \xi_1: &\ \lambda=\lambda,\notag \\
	{\rm coefficient\ of\ } \eta_1:&\ 0=0,\notag \\
	\label{2,0} {\rm coefficient\ of\ }\xi_1^2:& \left(\lambda^2-\lambda\right)a_2^2=G_{1,2}^2,\\
	\label{1,1} {\rm coefficient\ of\ }\xi_1\eta_1:& \left(1-\lambda\right)a_2^1=G_{1,2}^1,\\
	\label{0,2} {\rm coefficient\ of\ }\eta_1^2:& \left(\lambda^{-2}-\lambda\right)a_2^0=G_{1,2}^0,\\
	\label{3,0} {\rm coefficient\ of\ }\xi_1^3:& \left(\lambda^3-\lambda\right)a_3^3=G_{1,2}^1b_2^0+G_{1,3}^3+2a_2^2G_{1,2}^2,\\
	\label{2,1} {\rm coefficient\ of\ } \xi_1^2\eta_1:&\\
	  i\lambda\alpha_{2} + \left(\lambda-\lambda \right)a_{3}^{2}&= \left( a_{2}^{2} + b_{2}^{1} \right)G_{1,2}^{1} +2a_2^1G_{1,2}^2+2b_2^0G_{1,2}^0+ G_{1,3}^{2},\\
    \label{1,2} {\rm coefficient\ of\ }\xi_1\eta_1^2:&\\
     \left(\lambda^{-1}-\lambda\right)a_3^1&= \left( a_{2}^{1} + b_{2}^{2} \right)G_{1,2}^{1}+G_{1,3}^1+2a_2^0G_{1,2}^2+2b_2^1G_{1,2}^0,\\
	\label{0,3} {\rm coefficient\ of\ } \eta_1^3:& \left(\lambda^{-3}-\lambda\right)a_3^0=G_{1,2}^1a_2^0+G_{1,3}^0+2b_2^2G_{1,2}^0.
\end{align}
From this series of equations, we can see that from (\ref{2,0}) to (\ref{0,2}), we can directly solve the value of $a_2^0,a_2^1$ and $a_2^2$. Then plug in these values in (\ref{3,0}), (\ref{1,2}) and (\ref{0,3}) (and notice the conjugate relation (\ref{conj_relation})) , we can solve out $a_3^0,a_3^1$ and $a_3^3$ (notice that \(\lambda^{k} \neq 1\) for \(k = 1,2,3,4\) if
\(n\) is large enough). However, this process fails for $a_3^2$ since it has coefficient zero in (\ref{2,1}), and that is why $\alpha_2$ has to be involved in the normal form. We can then just let $a_3^2$ be zero, and find the exact value of $\alpha_2$ to satisfy the equation (\ref{2,1}). 

Now we find the value of $\alpha_2$.

From (\ref{2,1}),
\begin{align}\label{eq1}
\alpha_{2} &= -\frac{i}{\lambda}\left\lbrack \left( a_{2}^{2} + \overline{a_{2}^{1}} \right)G_{1,2}^{1}++2a_2^1G_{1,2}^2+2\overline{a_2^0}G_{1,2}^0  + G_{1,3}^{2} \right\rbrack,
\end{align}
and from (\ref{2,0}), (\ref{1,1}) and (\ref{0,2})
\begin{equation}\label{eq2}
 a_{2}^{2}=\frac{1}{\lambda( \lambda-1)}G_{1,2}^{2},\ \  a_{2}^{1}=\frac{1}{(1 - \lambda)}G_{1,2}^{1},\ \  a_{2}^{0}=\frac{1}{\lambda^{-2}-\lambda}G_{1,2}^{0}
\end{equation}
Combine (\ref{eq1}), and (\ref{eq2}), we have
\begin{equation}\label{alpha2}
\alpha_{2} = -i\left( \frac{G_{1,2}^{2}G_{1,2}^{1}}{\lambda^2( \lambda-1)} + \frac{\left|{G_{1,2}^{1}}\right|^{2}}{\lambda-1 } +\frac{2G_{1,2}^1G_{1,2}^2}{\lambda(1-\lambda)}+\frac{2\left|{G_{1,2}^{0}}\right|^{2}}{\lambda^3-1}+ \frac{G_{1,3}^{2}}{\lambda} \right).
\end{equation}

Moreover, we can see that $a_2^0,a_2^1$, $a_2^2$ are of order $\mathcal{O}(n^{-1})$ and $a_3^0,a_3^1$, $a_3^3$ are of order $\mathcal{O}(n^{-2})$. 
Therefore, we have the observation that $G_{i,j}^k$, $\phi_j$ and $\psi_j$ are of order $\mathcal{O}(n^{-j+1})$. Notice that we have checked that all polynomials of degree less than or equal to 3 coincide in equation \ref{Birk_eq}, then all remaining terms are polynomials of degree higher than 4 and the observation about $G_{i,j}^k$, $\phi_j$ and $\psi_j$ shows that they are all $\mathcal{O}_4\left(n^{-3}\right)$.
Thus, \ref{Birk_eq} is solved using this coordinate transform and the existence of conjugacy is proved.

Finally, we can see that $H_{2,2}$ is identity up to $\mathcal{O}(n^{-1})$ accuracy, thus $H_{2,2}(\widetilde{U})\sim \widetilde{U}$, up to $\mathcal{O}(n^{-1})$ accuracy. Since 
$$U=0.99P\left(	\frac{3}{2048},  -\frac{3}{2048}+\frac{3\sqrt{2}i}{1024}, -\frac{3}{1024},  \frac{3}{2048}-\frac{3\sqrt{2}i}{1042}\right),$$

$$V=0.98PP\left(	\frac{3}{2048},  -\frac{3}{2048}+\frac{3\sqrt{2}i}{1024}, -\frac{3}{1024},  \frac{3}{2048}-\frac{3\sqrt{2}i}{1042}\right),$$
we can conclude that if for $n$ large enough, $\widetilde{V}\subset H_{2,2}(\widetilde{U})$.

\end{proof}

\end{document}
